% ------------------------------------------------------------------ Abstract
\begin{simtext}
\noindent{\sffamily\bfseries\small
Pixel-level segmentation of microaneurysms, hemorrhages, hard exudates and soft exudates makes automated grading of diabetic retinopathy interpretable, and dozens of methods now report results on the same few public datasets. Whether those numbers can be compared is unclear: splits, input resolution and metric implementations differ, and many baselines are copied from earlier papers. Here we assessed \rn{n.models} model variants on \rn{n.sources} test sets from six datasets (\rn{n.images} images) under one protocol. Of \rn{funnel.dr} fundus-specific methods, \rn{funnel.train} came with training code and \rn{funnel.full} with everything needed to rerun them; a widely reused 11-threshold implementation of the area under the precision--recall curve scored a released checkpoint at \rn{e1.m2mrf.IDRiD.fn} where the exact value was \rn{e1.m2mrf.IDRiD.exact}. After unified retraining of \rn{n.main} models, \rn{cd.first.model} had the best average rank, differences among the top \rn{cd.ntop} models were within the critical difference, and performance outside the training source depended on the annotation style and field of view of the test images: it held on pixel-labelled standard-field sources and fell by a third to a half on coarsely annotated, region-labelled and ultra-widefield images. Annotation granularity, input resolution and lesion size explained more variance than architecture. We release code, splits and probability maps as a benchmark.}
\end{simtext}

\vspace{8pt}

% -------------------------------------------------------------- Introduction
\noindent Diabetic retinopathy (DR) is graded from lesions that are visible in colour fundus photographs: microaneurysms (MA), hemorrhages (HE), hard exudates (EX) and soft exudates or cotton wool spots (SE). A model that outlines these lesions pixel by pixel shows a grader why an image received its grade, supports screening programmes in which the decision must be auditable, and provides quantitative lesion burden for follow-up\cite{sebastian2023drlesionsurvey,tahir2024microaneurysmreview}. Public datasets with pixel-level labels, first IDRiD\cite{porwal2018idrid,porwal2020idridchallenge} and DDR\cite{li2019ddr} and later FGADR\cite{zhou2021fgadr}, Retinal-Lesions\cite{wei2020retinallesionslesionnet} and TJDR\cite{mao2023tjdr}, have turned the task into a standard benchmark, and methods built for it now span multi-scale convolutional networks\cite{guo2019lseg,liu2023m2mrf,xu2024hacdrnet,he2022pmcnet}, relation and transformer designs\cite{huang2022rtnet,he2023h2former}, semi-supervised schemes\cite{zhou2019collaborativesemisupervised,ullah2023ssmdunet} and adaptations of segmentation foundation models\cite{kirillov2023sam,ravi2024sam2,ma2024medsam,li2024tpdrseg}.

Several properties of the task make the reported progress hard to read. (1) The datasets are small and heterogeneous. IDRiD has 81 annotated images, its official test set 27, and soft exudates appear in about half of them; other sources were annotated with ellipses and polygons by many graders, or as regions with confidence levels by four experts\cite{kauppi2007diaretdb1}, so that ``ground truth'' means different things from one source to the next. (2) The lesions are small. A microaneurysm spans about half a percent of the field-of-view diameter, and shrinking an image to the 224-pixel input of a common backbone removes \rn{vanish.DDR.224.MA}\% of the microaneurysms and \rn{vanish.DDR.224.EX}\% of the hard exudates from the DDR training labels before any model is trained. (3) Evaluation is not standardized. Papers differ in the split (official or their own), the resolution at which predictions are scored, whether pixels outside the field of view count, and how the area under the precision--recall curve (AUPR) is computed; a widely reused implementation samples the curve at 11 thresholds. (4) Baseline numbers travel between papers without being rerun, sometimes across several generations and across different splits. (5) Claims about foundation models mostly come from the papers that propose an adaptation, with zero-shot use, prompting from ground truth and fine-tuning reported side by side.

Earlier syntheses do not resolve these issues. Surveys collect the numbers that papers report\cite{sebastian2023drlesionsurvey,tahir2024microaneurysmreview}. The closest study, by Playout and Cheriet\cite{playout2024crossdatasetgeneralization}, trained U-Net-type baselines across five datasets and showed that annotation style governs cross-dataset transfer; it did not rerun fundus-specific methods, did not audit published numbers and did not include foundation models. Outside ophthalmology, standardized reassessments have repeatedly found that part of the reported progress disappears once baselines are configured with care\cite{isensee2021nnunet,isensee2024revisited} and that rankings on small test sets are fragile\cite{maierhein2018rankings}. A recent assessment of T cell receptor--epitope predictors\cite{lu2026tcrepitopenaturemethods} supplies the design we follow: evaluate the original models, retrain everything that can be retrained under one protocol, and judge both on independent data.

We therefore assessed methods for DR lesion segmentation along five lines. (1) Systematic collection: a fixed literature search and a ledger of 86 records, \rn{n.sources} test sets with frozen splits, and \rn{n.models} registered model variants with a reproducibility level for each. (2) Audit of published numbers: which split, which metric implementation and which provenance stand behind each reported value. (3) Original models: released checkpoints scored with the authors' metric and with ours, and the authors' pipelines rerun from scratch. (4) Unified retraining: one protocol for data, training amount, augmentation, checkpoint selection and evaluation, three seeds, and five test sets from sources that were never used for training. (5) Factors: annotation granularity, training source, training amount, input resolution, lesion size, loss and the way a foundation model is used, each varied alone. The study answers three questions: how many published numbers can be reproduced and where the gaps come from (Q1); which family of methods is better once conditions are equal, and whether the answer survives on independent data (Q2); and which properties of the data and the protocol decide performance (Q3).
